\subsection{Which target quotients the descended columns must detect}
\label{rank:section}

Fix a scale and a translate in Section~\ref{geo:translate}. Suppress
the translate index. Write \(R\) for the generic fibre of the dual
target, \(d=\dim_K R>0\), \(r_i=\dim_K S_i\), and \(r=\sum_i r_i>0\).
The rational word maps \(B_i:S_i\to R\) retain their divisor shifts.
For \(W\subsetneq R\) defined over \(K\), let \(\pi_W:R\to R/W\),
\(\zeta_W=\dim(R/W)/d\), and set
\[
 q_W=\frac{1}{r}\sum_i\operatorname{rank}_K(\pi_W B_i),
 \qquad \tau=\min_{W\subsetneq R}q_W/\zeta_W .
\]
The minimum exists since the ranks are bounded integers. Taking
\(W=0\) gives \(0\le\tau\le1\). These tests use ranks over the number
field. The \src{src/Formalism/HarderNarasimhan/Statement.lean}{36}{original HN cutoffs} are submodules
over the ring of integers, so their generic fibres are eligible
\(W\)'s.

\begin{proposition}[A slope comparison from ranks on quotients]\label{rank:transfer}
Let \(s=\sum_i r_i(\mu(S_i)-h_i)/r\), and let \(U\) be the largest
original source slope after subtracting its allowance \(h_i\) and
the common centre \(s\). Write \(J=\sum_i g(B_i)/r\) for the cost in
Lemma~\ref{trmean:graph}. If \(\tau>0\), then
\[
                   s-\mu(R)\le\frac{J+(1-\tau)U}{\tau}.
\]
It suffices to have \(q_W\ge\tau\zeta_W\) on the proper target HN
cutoffs. For integral columns satisfying the prescribed heights
one may take \(J=0\).
\end{proposition}
\begin{proof}
Choose source and target slopes with their normalized rank weights
at this translate. Denote the centred repaired source by \(S'\)
and the centred target by \(T\). The source law has mean \(-J\)
and is bounded above by \(U\), by the integral inclusion into the
original source. At a real threshold \(v\), take the target cutoff
of slope at least \(s+v\). The integral map from the repaired word
source sends its cutoff at \(s+v+h_i\) into this target cutoff, by
the filtration inequality in Theorem~\ref{trmean:trace}. Therefore
\[
       \tau\Pr(T<v)\le q_W\le\Pr(S'<v)
\]
when the target cutoff is proper. For the whole target the left
side is zero. Integrate the resulting probability inequality over
\(-\infty<v<U\). The integral of each indicator is the positive
part of \(U\) minus the corresponding slope, hence
\[
 \tau(U-\mathbb ET)\le\tau\mathbb E(U-T)_+
        \le\mathbb E(U-S')_+=U+J.
\]
Rearrange and use \(\mathbb ET=\mu(R)-s\). The lemma on repairs
applies before taking ranks since the repaired source spans the
same generic fibre.
\end{proof}

There are finite certificates for the rank condition. A set of
words whose combined matrix has a nonzero \(d\)-rowed minor spans
\(R\). In every quotient its sum of word ranks is at least the
dimension of that quotient. If \(L\) disjoint sets of words have
such minors then \(\tau\ge dL/r\). More generally give spanning
sets \(C\) weights \(a_C\ge0\) with \(\sum_{C\ni i}a_C\le1\) for
each word. Summing the quotient inequalities gives
\(\tau\ge(d/r)\sum_C a_C\). Each minor can be checked at one place
of integral reduction. For a test restricted to a specified
codimension, one may instead verify the rank inequalities for every
quotient of that dimension over a single finite residue field at
which all columns are integral in lattice bases. Indeed, the
intersection of a rational subspace with the target DVR lattice
is saturated. Its reduction has the same codimension, and reduction
can only decrease the ranks of the induced maps into the free
quotient.

\begin{proposition}[The quotients covered by the original frame data]\label{rank:wide}
Use the full word maps \(A_i^{\,f}:V_i^{\,f}\to\mathcal H\) at the
selected infinite place, including the scalar \(\sqrt{\gamma_i}\).
The actual source isometries are \(V_i\hookrightarrow V_i^{\,f}\).
Projection to the \(d\)-dimensional target then gives the descended
\(B_i\) up to nonzero word scalars, in the coordinates of
\code{frame\_at\_of}. Put \(n=\dim\mathcal H\), \(r_f=\sum_i\dim V_i^{\,f}\),
\(\alpha=1-r/r_f\), \(\rho=d/n\), \(Q=\sum_i A_i^{\,f}A_i^{\,f*}\),
and \(M=\operatorname{tr}Q\). Suppose \(M\ge r_f/2\), and write
\[
 \eta=\frac{n\operatorname{tr}(Q^2)}{M^2}-1,\qquad
 e=\frac1{r_f}\sum_i
      \|A_i^{\,f*}A_i^{\,f}-I\|_{\rm HS}^2,\qquad
 D_f=\eta+5\alpha+4e .
\]
Every quotient with \(\rho\zeta_W\ge2D_f\) satisfies
\(q_W\ge(\rho/10)\zeta_W\).
\end{proposition}
\begin{proof}
Let \(P\) project in \(\mathcal H\) onto the orthogonal complement
of \(W_{\mathbf C}\) inside the complex target. Its normalized trace
is \(z=\rho\zeta_W\). In the \(i\)-th full source \(P A_i^{\,f}\)
vanishes on the embedded kernel of \(\pi_W B_i\). The fraction of
full source dimensions outside these kernels is exactly
\(b=\alpha+(1-\alpha)q_W\). Positivity and the Hilbert--Schmidt
Cauchy--Schwarz inequality give
\[
 z-\sqrt{z\eta}\ \le\ \frac{\operatorname{tr}(PQ)}{M}
                  \ \le\ 2(b+\sqrt{be}).
\]
On the left pair \(P\) with \(nQ/M-I\) in normalized trace.
On the right use an orthonormal basis of the kernel complement
in each full source. Bounding the squared norms there uses \(I\)
and \(A_i^{\,f*}A_i^{\,f}-I\), and \(r_f/M\le2\).
Now \(2\sqrt{z\eta}\le z+\eta\) and \(4\sqrt{be}\le b+4e\).
It follows that \(z\le\eta+5b+4e\le D_f+5q_W\). Under the stated
size condition this gives \(z\le10q_W\).
\end{proof}

For comparison with the source, the
\src{src/Formalism/SharpRestriction/GlobalSchurBounds/TopDensityNGeSixteen.lean}{538}{original frame comparison}
passes the full
maps, isometries, weights and actual \code{Bk} to the joint lattice
theorem. The latter obtains trace and diagonal estimates for
compressed maps in \code{ramification\_linear\_case\_holds}.
Proposition~\ref{rank:wide} tests the full maps directly. For the
uncompressed trace the same source bound specializes to
\(\eta\le5\sigma\) by using identity isometries and zero rank loss
under its original size and Gram hypotheses. A diagonal defect
\(e\le1/4\) gives \(M\ge r_f/2\) by Cauchy--Schwarz.

Thus only the integer codimensions \(1\le b<2nD_f\) need an
additional rank test if one seeks a constant at most \(\rho/10\).
At \(2nD_f\le1\) this list is empty. In the original period
construction the period image surjection in \code{IndefUse.lean},
line 463, identifies the space of rows. The
\src{src/Formalism/LocalToGlobal/IndefiniteForms/IndefSide.lean}{388}{columns on these rows}
are pairings with the actual word images. A common missed
row is a row orthogonal to all those images, in the coordinates of
the common row isometry. The deeper realization in \code{RowSideTie.lean},
line 253, supplies a bijection of word labels. The kernel package in
\code{RecBridgeSel.lean}, lines 272--280, identifies the period at
each original label. A nonzero column at each word would say that
each word pairs with some row. For a one-dimensional residual quotient
we need a word pairing with its particular dual row.

\begin{proposition}[A missed line under the pairwise Gram bounds]\label{rank:blind}
For every odd prime \(p\ge41\) and every integer \(k\ge1\) there
are \(p(p+1)\) full word maps of rank \(k\) in dimension \(n=kp+1\)
satisfying the three bounds of \code{InputSG}. They meet the word
size and count comparisons with \(C_{\rm pr}=1\), \(\Gamma^2=p\)
and \(\sigma=2/p\). Every word vanishes in the same one-dimensional
rational quotient.
\end{proposition}
\begin{proof}
In \(\mathbf C^p\) take the coordinate basis and, for each
\(a\in\mathbf F_p\), the \(p\) unit vectors
\[
 v_{a,b}(x)=p^{-1/2}\exp(2\pi i(ax^2+bx)/p),
                       \qquad b,x\in\mathbf F_p .
\]
Each fixed \(a\) gives an orthonormal basis. Two vectors in different
bases have squared inner product \(1/p\). For two quadratic bases
this follows by squaring the absolute value of the unnormalized
sum and writing \(x=x'+h\). For \(h\ne0\) the inner sum over \(x'\)
is a nontrivial additive character sum and is zero. Only \(h=0\)
contributes \(p\). This is the prime-dimensional construction of
\href{https://doi.org/10.1016/0003-4916(89)90322-9}{Wootters and Fields}.

For each resulting unit vector \(v\), use the isometry
\(z\mapsto(v\otimes z,0)\) into
\((\mathbf C^p\otimes\mathbf C^k)\oplus\mathbf C\).
Its unnormalized columns are defined over the \(p\)-th cyclotomic
field, with the factors \(p^{-1/2}\) placed in \(\sqrt{\gamma_i}\).
The diagonal Gram defect is zero, and each off-diagonal squared
Hilbert--Schmidt norm is zero or \(k/p=k^2/D_*\) with \(D_*=kp\).
Take \(e_0=\eta_0=0\), \(\eta_1=1/(kp)\). Then
\(n=(1+\eta_1)D_*\), \(\eta_1\le\sigma/2\),
\(k\Gamma^2/D_*=1\), and the word count exceeds
\(\Gamma^2/\sigma\). There is no compression. Summing the projections
over the \(p+1\) bases gives \(Q=(p+1)I_{kp}\oplus0\), so
\(\eta=1/(kp)\). The last coordinate receives zero from every word.
Giving this line slope \(-L\) in an orthogonal direct sum with
standard visible coordinates leaves the frame numbers fixed and
decreases the target mean by \(L/n\).

The missed line also persists with a nontrivial ramified rank
ratio. Take \(k=m(T-1)^f\) and \(k'=m(T-2)^f\) with \(2k'\ge k\),
and restrict each full map to \(k'\) coordinates of its source.
Use the first \(k'\) coordinates throughout one of the bases and
the last \(k'\) throughout a second basis. These two groups of
words already span \(\mathbf C^p\otimes\mathbf C^k\). The other
words may use either choice. The actual rank is \(k'\) throughout,
all full Gram tests still hold, and exactly the last row is missed.
These examples test the stated linear and size comparisons. The
descended periods on an arithmetic host have further structure.
\end{proof}

\paragraph{Size required for the window.}
Return to the translates and take \(y=x^\varepsilon\) and
\(P_0=n_Ta_E(1+m^2/k)+m(B_0+N_c)\). Assume the lower comparison
\eqref{trmean:lower}, the original level and variation bounds,
integral columns with their height allowances, and
\(\mathbb E(\mathcal A_u-G_u)\le C P_0\) for a uniform \(C\),
\(c_s\mathbb E I_u\ll P_0\). Under the size hypotheses in
Corollary~\ref{trmean:reserve}, including \(e_{\rm div}\le2\), we have
\[
 K_0\ll y^2P_0,\qquad
 \mathbb E U_u\le K_3+K_7+c_s\mathbb E I_u\ll yP_0 .
\]
The first bound uses the separate terms of \(K_0\) in the proof
of that corollary. On a given scale suppose
\(\min_u\tau_u\ge c\,y^{-4+\xi}\), with fixed \(c>0\) and
\(0<\xi<4\) on the family. Proposition~\ref{rank:transfer} and
\eqref{trmean:lower} give
\[
          \mathbb E D_u\le C'(y^2+y^{5-\xi})P_0
\]
for a uniform \(C'\). For each fixed \(\varepsilon>0\) the right side is at most \(y^5P_0\)
beyond one uniform threshold. The bounds for \(B_0,N_c\) in
Corollary~\ref{trmean:reserve} then give \eqref{scope:average}.
Thus even the rank fraction \(c\,y^{-3}\) suffices under the listed
arithmetic inputs. If \(\rho\ge1/2\), the large quotients already
pass this test eventually by Proposition~\ref{rank:wide}.
No bound on the minimum target slope enters this implication.

\begin{corollary}[One minor for the remaining HN quotients]\label{rank:one-minor}
Use the original per-word lower and variation inequalities in the
proof of Proposition~\ref{trmean:assembly} and
Corollary~\ref{trmean:loss}, and the other arithmetic size, column
height and loss bounds just stated. Suppose \(\rho\ge1/2\) at each
translate. Among the proper target HN cutoffs with
\(0<\dim(R/W)<2nD_f\), take the smallest subspace \(W_*\), if there
is one. It suffices that on each such quotient \(R/W_*\) some set
of words \(C\) has surjective combined image and
\[
                r_C:=\sum_{i\in C}r_i\le C_0\,d\,y^{4-\xi}
\]
for fixed \(C_0>0\), \(0<\xi<4\). Then \eqref{scope:average} holds
on the stated scales beyond one uniform threshold. Surjectivity
here can be witnessed by a single nonzero minor of size
\(\dim(R/W_*)\) in the projected columns.
\end{corollary}
\begin{proof}
The target HN cutoffs form a finite flag, so all the small quotients
factor through \(R/W_*\). At this translate assign word \(i\) the
probability
\[
                \omega_i=\frac{r_i}{2r}
                         +\mathbf1_{i\in C}\frac{r_i}{2r_C}
\]
and distribute it uniformly on the slopes of the word. In a quotient
with relative dimension \(\zeta_W\) the weighted fraction of
detected source ranks is
\(\sum_i\omega_i\operatorname{rank}(\pi_W B_i)/r_i\).
On the large quotients the first half gives at least
\((\rho/20)\zeta_W\) by Proposition~\ref{rank:wide}. On the small HN
quotients the second half gives at least \((d/(2r_C))\zeta_W\).
If there is no small HN quotient use the original word weights.
Thus the weighted HN cutoffs have a common lower ratio
\(\gg y^{-4+\xi}\). Recentring the source changes the numerical
thresholds but preserves the entire target flag.

The per-word estimate
\(G_u-\mathcal M_u\le\mu(S_i)-h_i-\mu_u+K_0+\ell_i\)
can be averaged with these probabilities. The pairwise inequality
\(\lambda_i^{\max}-h_i-(\mu(S_j)-h_j)\le K_3+\ell_j\)
gives \(U^\omega_u\le K_3+\sum_i\omega_i\ell_i\), and
\(\ell_i\le c_sI_u+K_7\) for each word. The proof of
Proposition~\ref{rank:transfer} applies to these rank probabilities,
with zero repair cost. The preceding numerical absorption proves
the assertion with the original \(D_u\) and \(E_x\). The choice of
probabilities is internal to the slope estimate.
\end{proof}

Write \(b_*=\dim(R/W_*)\). If the combined image spans the quotient,
one can select at most \(b_*\) words by enlarging the image one word
at a time. The original \code{InputSG} even bounds the cost of each
word. Its per-word diagonal estimate implies
\(\dim\ker A_i^{\,f}\le e_0\dim V_i^{\,f}\), since
\(A_i^{\,f*}A_i^{\,f}-I\) acts by \(-1\) on this kernel. Thus
\(r_i/d\le1/(\rho(1-e_0))\le8/3\) for the stated target ratio and
\(e_0\le1/4\). When \(b_*=1\), one nonzero projection alone suffices.

The original word size comparison gives a sharper budget for higher
codimensions. For \(k_i=\dim V_i^{\,f}\) the condition
\code{ChProviso\_at} gives \(k_i\le C_{\rm pr}D_*/\Gamma^2\).
In the original full Gram hypotheses
\(D_*/r_f\le\sigma\le1/20\) and
\(\sqrt{e_0},\eta_0\le\sigma/2\). The trace inequality
\(M^2\le n\operatorname{tr}(Q^2)\) implies
\[
 \frac{D_*}{n}\le
 \frac{1+\eta_0+(D_*/r_f)(1+\sqrt{e_0})^2}
      {(1-\sqrt{e_0})^2}<2 .
\]
Indeed \(M\ge(1-\sqrt{e_0})r_f\). The diagonal terms of
\(\operatorname{tr}(Q^2)\) are at most
\((1+\sqrt{e_0})^2r_f\) by the Hilbert--Schmidt triangle inequality,
and the off-diagonal terms sum to at most
\((1+\eta_0)r_f^2/D_*\) by \code{InputSG}. Consequently the
selected words cost \(r_C/d\le4C_{\rm pr}b_*/\Gamma^2\).
In particular \(b_*\ll\Gamma^2y^{4-\xi}\) and surjectivity onto
this single quotient suffice. The original source obtains the
full-word ranks and their operator estimates in
\code{FrameReal\_at}, with the
\src{src/Formalism/SizeWindowBrackets/FrameTraceRatioFinEquiv.lean}{30}{pairwise Gram bounds}
and the \src{src/Formalism/SharpRestriction/GlobalSchurBounds/Chosen.lean}{53}{two-sided word size}.
A uniform certificate on the actual
residual quotient is still required.

The ramified rank formula explains why an estimate for the fraction
\(\alpha\) alone cannot dispose of the issue. Assume the original
identification \(r_i=k_i(1-1/(T_i-1))^f\) with \(T_i>2\).
The consecutive integers \(T_i-1,T_i-2\) are coprime. Clearing
denominators shows that for some positive integer \(m_i\)
\[
        k_i=m_i(T_i-1)^f,\qquad r_i=m_i(T_i-2)^f.
\]
Let \(N\) be the number of words. Since
\(r_f\le NC_{\rm pr}D_*/\Gamma^2\), the same trace bound gives
\[
 \frac{n\alpha}{\Gamma^2}
 \ \ge\ \frac{1}{2C_{\rm pr}N}\sum_i(k_i-r_i)
 \ \ge\ \frac{\min_i\big((T_i-1)^f-(T_i-2)^f\big)}{2C_{\rm pr}} .
\]
For \(f\ge2\) the last difference is at least
\(f(\min_i T_i-2)^{f-1}\), by the binomial theorem.
On the original scales \(\min_i T_i\ge b=\exp(y)\). Thus the
upper cutoff \(2nD_f\), measured in units of \(\Gamma^2\), already
exceeds every fixed power of \(y\) from its compression term alone.
This statement bounds the cutoff used in Proposition~\ref{rank:wide}.
It gives no lower bound for the size of an actual missed HN quotient.

There is a direct arithmetic formulation of the minor itself.
Work in the original row coordinates at \(\mu_0\) and put
\[
 \begin{split}
 U_{\rm sp}&=V_0\cap\ker\iota,\\
 K_{\rm per}&=\{F\in V_0:\ 
    \operatorname{perCol}(i,\mu_0,\phi,F)=0
                       \ \text{for every actual word }i,\phi\}.
 \end{split}
\]
The row map \(\iota:V_0\to EE\) is onto under the original row
realization. The identity \(\operatorname{perCol}(i,\mu_0,\phi,F)
=B_i(\phi)(\iota F)\) for \(F\in V_0\) identifies the common row
annihilator with \(\iota(K_{\rm per})\), hence with
\(K_{\rm per}/U_{\rm sp}\). The original
\code{periodIdentity\_iff\_period\_kills} in \code{DataBu.lean},
line 49, expresses factorization by the inclusion
\(U_{\rm sp}\subseteq K_{\rm per}\). The word domain here is the
ramified one. In \code{RhoDetPrimeEqMuTildeMulChiPow.lean}, lines
165--191, \code{eMu} identifies that space and \code{jN\_basis}
shows that images of an actual rational word basis span it over
\(\mathbf C\). Testing all such images or their complex span gives
the same annihilator, by linearity in \(\phi\).

Put \(A_*=(R/W_*)^*_{\mathbf C}\), viewed in the row coordinates
dual to the column target. The precise nonvanishing needed for
the full set of words is
\[
                K_{\rm per}\cap\iota^{-1}(A_*)=U_{\rm sp}.
\]
Indeed it says that the common annihilator meets \(A_*\) in zero,
which by finite-dimensional duality is equivalent to surjectivity
onto \(R/W_*\). The companion \code{PeriodRadical.lean} proves this
annihilator criterion with the actual number-field word vectors in
\code{jN}, taking the degree-zero factorization and row surjectivity
as inputs. For rank one this asks for one nonzero period on the
prescribed line. In \code{CayleySignature.lean}, line 138, the row,
vector and test witnessing a nonzero automorphic period may be
chosen together. That quantifier does not prescribe this line.

The row-lift identity in \code{PrincipalHermitian\allowbreak Lattice.lean},
lines 101--105, diagonalizes \(\iota_Q\iota_Q^*\). For the selected
word images let \(P_*\) be the orthogonal
projection onto the directions of the residual quotient at the chosen
infinite place. Here the operator whose invertibility would give the
minor is
\[
   \sum_i (P_* B_i)(P_* B_i)^* .
\]
Its quadratic form at \(a\) is \(\sum_i\|(P_*B_i)^*a\|^2\), so its
kernel is exactly the common annihilator in those directions, using
the Hermitian identification with the dual. In \code{CxHecke.lean},
lines 145--149, the row components are simultaneous eigenspaces.
The multiplicity \code{nRows} in \code{DataRows.lean}, lines 209--215,
counts all rows of a given label. Even in a Hecke-stable kernel,
one detected vector in a two-dimensional space of the same eigenvalues
can leave a one-dimensional kernel: every Hecke operator acts by a
scalar on this whole space. The stronger period statement
\code{LiftParameterConditions} in \code{CayleySignature.lean}, lines
163--166, prescribes the representation on the column side and still
chooses the row witnessing nonvanishing.

\paragraph{A period test at the chosen level prime.}
Fix one translate, a nonzero vector \(a_*\) on the prescribed dual
line, and a lift \(F\in V_0\) with \(\iota F=a_*\). Keep this lift
fixed for all words. The source first chooses the code, weights and
full dimension data. For word \(i\) it then supplies a norm prime
in \([T_i^*,(1+1/(fk))T_i^*]\), before constructing the final row
lattice. The selection theorem does not impose pairwise distinctness
on these primes. On one realized family the word-specific map
\(e_{i,0}:\mathrm{Mt}_i\longrightarrow\mathrm{Amb}_{\mu_0}\)
identifies the row used in its full period:
\[
 L_{i,F}(\phi)=
  \mathrm{period}_i\bigl(\mathrm{lowerMap}_i
       (\mathrm{evz}_i(\mathrm{col}_i\phi))\bigr)(e_{i,0}^{-1}F).
\]
The column chain and this row identification both depend on \(i\).
\src{src/Development/AbelianVarietyTheory/Proofs/GlobalInterpolation/ExistsObjectsOfAdmissibleTwistDatum.lean}{97}{order of the level choice}
and \src{src/Development/AbelianVarietyTheory/Proofs/GlobalInterpolation/ExistsLevChOfLevelPrime.lean}{92}{the individual intervals};
\src{src/Formalism/LocalRingCardinality/RhoDetPrimeEqMuTildeMulChiPow.lean}{165}{word-specific row and source identifications}
and \src{src/Formalism/LocalRingCardinality/DataKernels.lean}{192}{the full period}.

In the tower argument the upper representation is fixed before the
upper column and row are chosen together. At a descent step the
column is expanded in a finite theta sum. Nonvanishing of the sum
selects a summand while retaining the higher row. Unfolding chooses
a new Schwartz function, and its theta lift of that higher row is
the lower row. The terminal witness thus contains an existential
position and a row. In the actual use of the unfolding identity
the proof obtains a primed tuple of witnesses and extracts a nonzero
summand to prove nonvanishing of an archimedean functional.
The companion \code{PeriodWitnessChain.lean} reconstructs the
descent from the six stated tower inputs. It keeps a supplied
higher row whenever a nonzero test on that row is supplied.
\src{src/Formalism/SharpRestriction/CuspSpanNormBounds/CayleySignature.lean}{154}{theta expansion, upper witness and descending test},
\src{src/Formalism/SharpRestriction/CuspSpanNormBounds/ThresholdCubicInequalitySignPattern.lean}{1671}{the row in the descending induction} and
\src{src/Formalism/SharpRestriction/Divisibility/InnerProductFormulaField.lean}{390}{extraction of the archimedean summand}.

Changing the selected primes can change the target lattice because
the host construction uses the image of the word-to-prime map.
Comparing a prescribed HN line
across two such choices would require an identification carrying
that line and its period pairings across. Within a single realization
the fixed lift above gives the same row for all the word tests.
\src{src/Development/AbelianVarietyTheory/Proofs/GlobalInterpolation/ObjOnNarrowOfObjectsExist.lean}{720}{the marked host in the realization}.

There is substantial choice of the primes before this realization.
Fix any rational prime \(\ell\). Write \(U=T_i^*\), \(V=fk\),
\(M=N(g_0g_1)^2\) for the original norm modulus, and \(d_K\)
for the quadratic parameter denoted \(d\) in the norm-prime count.
The short-interval proof first establishes a lower bound for the
number of allowed primes, then retains just one of them. Write
\(A,c,c'>0\) for its absolute constants. Put
\[
 a_i=\left\lfloor\frac{c\sqrt{\log U}}{4\log\ell}\right\rfloor,
 \qquad M'=M\ell^{a_i}.
\]
For each fixed \(0<\eta<1\), every sufficiently late word window
\((U,U(1+1/V)]\) contains at least \(U^{1-\eta}\) admissible
level primes satisfying \(T_i\equiv1\bmod\ell^{a_i}\).
Indeed the source hypotheses give
\(\log N\ll x^{\varepsilon/5}\log^2 x\), \(d_K\mid N\),
\(\log g_j\ll x^\varepsilon\), \(V\le x^{12}\), and
\(\log U\gg x^{1/3+\varepsilon}\). Thus
\(\log(4d_KM^2)=o(\sqrt{\log U})\) and
\(\log V=o(\sqrt{\log U})\), uniformly in the word.
In particular \(\log(4d_K(M')^2)\le c\sqrt{\log U}\)
and \(1/V\ge\exp(-c\sqrt{\log U})\). The original count with
modulus \(M'\) gives the lower bound
\[
        \frac{c'U}{V(4d_K(M')^2)^A\log U}
        \ \ge\ U^{1-\eta}
\]
eventually, by taking logarithms. A norm generator congruent to
one modulo \(M'\) also supplies the required generator modulo \(M\),
and its norm is one modulo \(\ell^{a_i}\). The original excluded
primes are below the left endpoint at these scales. The per-word
level constructor accepts the chosen generators. The original
realization hypothesis is then applied to this very level choice
with its fixed \(d_i,D_*,\Gamma\). In particular any
fixed pair of word labels can receive distinct primes, even if
their windows coincide. For any fixed \(B\), the same greedy choice
works for at most \(x^B\) designated words. This quantifies prime
supply before the row lattice is formed. The companion
\code{LevelPrimeOrbitSupply.lean} uses the original short-norm
theorem to certify the finite count and avoidance of a prescribed
set of smaller cardinality.
\src{src/Development/AutomorphicArithmetic/ShortNorms/ShortNorm.lean}{8}{the uniform ray-norm count},
\src{src/Development/AutomorphicArithmetic/ShortNorms/levelPrimeN.lean}{36}{the level-window parameters},
\src{src/Development/AbelianVarietyTheory/Proofs/GlobalInterpolation/NormOnePrimeElements.lean}{15}{the norm generator and its congruence} and
\src{src/Development/AbelianVarietyTheory/Proofs/IwasawaDecomposition/ModEqOneOfNorm.lean}{47}{assembly with the selected generators}.

Use the original free finite model of the full level-zero columns.
Each orbit is a copy of \(H^{\mathcal V}\), where
\(H=\mathbf F_{T_i}^{\times}\), \(q_i=|H|=T_i-1\), and
\(|\mathcal V|=f\). The model may have many orbits. Its ramified
subspace consists exactly of functions with zero sum along each
coordinate fibre. This follows by transporting each local norm
operator through the model equivalence. The original source
constructs the model from the split level and its diamond operators.
\src{src/Formalism/LocalToGlobal/GlobalTheory/GL6/RamifiedRankEqFinrank.lean}{43}{transport of all norm kernels} and
\src{src/Formalism/LocalToGlobal/GlobalTheory/GL6/RamifiedLatticeOfRankSlopeLength.lean}{263}{the free model at the split prime}.

On one orbit put \(c(g)=L_{i,F}(\delta_g)\), using the delta
column transported back to the full space. For any \(a\in H^{\mathcal V}\)
use the integral-valued function supported on this orbit with values
\[
       w_a(g)=\prod_{v\in\mathcal V}
          (1-q_i\mathbf1_{\{g_v=a_v\}}).
\]
Every coordinate sum is zero. Its period is exactly
\(\sum_g w_a(g)c(g)\). The \(w_a\) span the ramified part on the
orbit over \(\mathbf C\): in each coordinate their factors are the
columns of the all-ones matrix minus \(q_i I\), which has image the
augmentation subspace. Tensoring these matrices gives the simultaneous
zero-sum space. Thus all these evaluations, on all orbits, vanish
exactly when word \(i\) misses \(a_*\).

There is a shorter exact test on each orbit. Fix a base point
\(b\in H^{\mathcal V}\). In coordinate \(v\) the vectors
\(\delta_{a_v}-\delta_{b_v}\), for \(a_v\ne b_v\), are a basis of
the zero-sum space. Their tensors form a basis of the ramified part.
Consequently one actual basis column of word \(i\) detects the line
if and only if, on some orbit and for some such \(a\),
\[
   \sum_{S\subseteq\mathcal V}(-1)^{|S|}c(g_S)\ne0,
   \qquad (g_S)_v=
     \begin{cases} b_v&v\in S,\\ a_v&v\notin S.\end{cases}
\]
Each individual test uses only \(2^f\) full periods at the same
level. A nonzero sum certifies detection immediately. To certify
vanishing this way one must exhaust the orbit bases.

There is an integral congruence even for arbitrary complex periods.
On the chosen orbit let
\(\Gamma=\sum_g\mathbf Z\,c(g)\subset\mathbf C\).
Since every \(w_a(g)\equiv1\bmod q_i\),
\[
     \sum_g w_a(g)c(g)\equiv\sum_g c(g)\pmod{q_i\Gamma}.
\]
If the class of the sum in \(\Gamma/q_i\Gamma\) is nonzero,
every anchor gives a ramified period detecting the line. In terms
of the integer relations among the orbit values, the test succeeds
exactly when no relation has all its coefficients congruent to
one modulo \(q_i\). It applies without an algebraicity hypothesis.

For a valuation test, suppose the coefficients on one orbit have a common
complex scaling after which they belong to the fraction field of a
discrete valuation ring. For a finite prime \(\mathfrak l\) of
residue characteristic dividing \(T_i-1\), it suffices to show
\[
 v_{\mathfrak l}\!\left(\sum_g c(g)\right)
       =\min_g v_{\mathfrak l}(c(g))<\infty
\]
in these coordinates. Indeed, divide by a coefficient of minimal
valuation. The normalized sum is nonzero modulo \(\mathfrak l\), and
every \(w_a(g)\) reduces to one. Hence every \(a\) gives a nonzero
ramified period. On the original word space, the images under
\(\mathrm{jN}_i\) of the number-field basis span over \(\mathbf C\).
Linearity and the descent identity give an actual basis column
nonzero on \(a_*\). The companion \code{RamifiedOrbitPeriod.lean}
checks this implication with the original \(\mathrm{PtilCol}\),
\(\mathrm{eMt}\), norm kernels and number-field basis. If the
lift is changed by a vector in \(U_{\rm sp}\), its ramified
evaluations stay fixed. Its two full orbit sums therefore also
agree modulo \(\mathfrak l\) whenever both tables are integral in
the same normalization.

The strengthened choice of level prime above gives a weaker
valuation requirement. Normalize \(v_{\mathfrak l}(\ell)=1\).
On one orbit of the resulting actual word, assume its period values
have a common algebraic scaling, write \(m_c=\min_g v_{\mathfrak l}(c(g))\)
in those coordinates, and suppose \(m_c<\infty\). It suffices that
\[
       v_{\mathfrak l}\!\left(\sum_g c(g)\right)-m_c<a_i .
\]
Indeed the difference between the weighted and unweighted sums
has valuation at least \(m_c+a_i\), as every \(w_a(g)-1\) is
divisible by \(\ell^{a_i}\). The weighted sum then has the same
finite valuation as the unweighted sum. Thus even a normalized
trace divisible by \(\ell\) can certify an actual detecting column.
It would suffice to find an orbit with finite normalized trace
valuation \(o(\sqrt{\log T_i^*})\) on one selected word per
required line. The estimate must hold for the row of the resulting
realization. Prime abundance alone gives no bound on this trace.

The same condition has an intrinsic form for arbitrary complex
periods. Define \(t_\ell(S)\) for \(S=\sum_g c(g)\ne0\) as the
largest \(k\) with \(S\in\ell^k\Gamma\), and set \(t_\ell(0)=\infty\).
The group \(\Gamma\) is finitely generated and torsion-free, hence
free. In any integral basis \(t_\ell(S)\) is the smallest
\(\ell\)-adic order of its coordinates. Detection follows whenever
\(t_\ell(S)<v_\ell(T_i-1)\), in particular when \(t_\ell(S)<a_i\).
For example a compatible homomorphism to a fixed finite-rank
integral lattice taking these traces to a fixed nonzero vector
would bound all the \(t_\ell(S)\) from above. Its construction
would have to respect the prescribed rows. For an orbit function
equal to one on a full coordinate fibre and zero elsewhere,
\(\Gamma=\mathbf Z\) and \(S=T_i-1\). The threshold is then
attained and every ramified test vanishes.

Two words at different odd level primes admit a common residue test
at \(2\). For example suppose the specified line lies in a fixed
two-dimensional row space on this realization. Lift a basis of
that space to \(F_1,F_2\). Choose one orbit at each of two words
\(i_1,i_2\) and form the \(2\)-by-\(2\) matrix with entries
\[
    C_{\nu j}=\sum_g L_{i_\nu,F_j}(\delta_g).
\]
If the delta coefficients are integral at a common place over \(2\)
after nonzero row and column scalings of the tables, and the scaled
\(C\) has nonzero determinant there, one original number-field
basis column from each word restricts independently to the plane.
Indeed every \(w_a\) reduces to one at both words. The resulting
two ramified columns have a nonzero complex determinant on the two
rows. Expand both columns in their respective descended bases and
use bilinearity of the determinant. These words then detect every
line in that plane between them. The same argument works directly
with two nonproportional vectors of mixed differences over
\(\mathbf C\). The choice of distinct primes by itself gives no
value for either determinant.

At the late levels \(T_i>2\). For \(c(a)\ne0\) write \(d(a,g)\)
for the number of differing coordinates. A sufficient archimedean test is
\[
         \sum_{g\ne a}\frac{|c(g)|}{(T_i-2)^{d(a,g)}}<|c(a)|.
\]
This follows by dividing the evaluation by \(w_a(a)\) and taking
absolute values. In particular any nonzero orbit table supported
on fewer than \(T_i-1\) entries detects the row, by taking a
coefficient of largest modulus. The coordinate-fibre example
above attains this support threshold.
The same vanishing holds for any table constant in one coordinate.
The remaining arithmetic question is whether at least one selected
word has a surviving ramified evaluation for the specified row.
The source's choices of \(T_i\) provide neither the residue of the
normalized full sum nor the concentration of its coefficients.

\paragraph{A bound when the lowest target slopes are available.}
Alternatively take \(H>0\), assume \(\rho_u\ge\rho_0>0\) everywhere
and retain the same full frame data at each translate. Let
\(T_{\rm out}=\mathbb E(-H-T)_+\), where the expectation includes
the target rank measure, and \(J_u\) denotes the repair cost at \(u\)
in Proposition~\ref{rank:transfer}. Applying the proof of
Proposition~\ref{rank:wide} to the target HN cutoff at \(-t\)
gives the following refinement of \eqref{trmean:budget-eq}:
\[
 \mathbb E(s_u-\mu_u)\le T_{\rm out}
  +\frac{H\mathbb E D_{f,u}
       +5\mathbb E\bigl(J_u+HU_u/(H+U_u)\bigr)}{\rho_0}.
\]
Indeed \(q_W\le\Pr_u(S'<-t)\) at that cutoff, including the zero
target codimension case by taking \(q_W=0\). Integrate
\(\rho_u\Pr_u(T<-t)\le D_{f,u}+5\Pr_u(S'<-t)\) from zero to \(H\).
The repair increases the clipped negative source mean by at most
\(J_u\). For any original centred slope \(a\le U_u\),
\[
       \min(H,(-a)_+)\le\frac{H(U_u-a)}{H+U_u}.
\]
This follows on the three intervals cut by \(-H,0\). Take the
source mean using \(\mathbb E_u a=0\), then the translate mean.
The rank and diagonal defects in the coefficient of \(H\) are
linear. The original size assumptions controlling the full
covariance may themselves involve square roots.

The companion \code{RankTailTransfer.lean} proves the quotient
rank bound for an individual cutoff, the rank test for a spanning
set of words, the passage from the finite cutoff inequalities to the mean and the scalar
estimates for the full frame and clipped source. The projection,
the arithmetic identifications and the window deductions above
are paper arguments. A lower bound on the ranks in the remaining
small quotients of the actual family, or an arithmetic bound on
their exceptional slopes, is needed to use this comparison on
arbitrarily late windows.
